%% GENERATED by tools/gen_latex.py -- edit content/hendri-agentic.py instead.
%% Build:  latexmk -pdf hendri-agentic.tex     (or: pdflatex hendri-agentic.tex, twice)
\documentclass[aspectratio=169,11pt,t]{beamer}
\usepackage{itbpro}

\coursetitle{Agentic AI}
\coursesubtitle{What you build around a language model --- and the more useful question of when you should not.}
\coursekicker{Module}
\coursesource{Module material for Designing and Building AI Products and Services}
\courseduration{3 hours (2 sessions)}
\coursepresenter{Hendri Karisma, M.T.}{Teaching Assistant}
\coursebrand{AI for Professional \textperiodcentered\ ITB}

\title{Agentic AI}
\author{Hendri Karisma, M.T.}
\date{}

\begin{document}
\covertitle

\framekicker{Learning outcomes}
\begin{frame}[shrink=25]{Session Objectives}
\leadpar{By the end of this session, participants can:}
\begin{isteps}
\item Distinguish \textbf{AI, AI workflow, AI agent, agentic system, and multi-agent system} using one operational test.
\item Describe the \textbf{agent loop} --- plan, act, observe, check --- and name what terminates it.
\item Design a \textbf{tool} properly: schema, description, validation, and the fact that only your code touches a real system.
\item Choose between the \textbf{four kinds of memory}, and say which problem each one solves.
\item Lay out a \textbf{tech stack} across model, orchestration, tools, memory, guardrails, and observability.
\item Apply a \textbf{three-question test} before splitting a system into multiple agents.
\item Build an \textbf{evaluation harness} for a non-deterministic system, and explain why traces come first.
\item Place a proposed deployment on the \textbf{autonomy ladder} and justify the level chosen.
\end{isteps}

\vskip 6pt
\vskip 4pt
\reslink{Site}{Course home}{../../index.html}
\reslink{Repo}{ai-agentic-demo --- single-agent and multi-agent cases}{https://github.com/situkangsayur/ai-agentic-demo}
\vskip 2pt
\end{frame}

\coursesection{01}{What the words mean}{Five terms that are used interchangeably and should not be.}

\framekicker{Setting up}
\begin{frame}[shrink=25]{This module sits on top of chapter 16}
{\fontsize{9.4}{12.6}\selectfont\color{fg2}Chapter 16 gave us a language model: a system that, given a sequence, produces a plausible continuation. It also gave us the honest limits --- hallucination, prompt sensitivity, no adaptation to genuine novelty.\par}\vskip 4pt
\leadpar{This module asks a different question: \textbf{given such a model, what do you build around it} so that it can affect the world rather than only describe it?}
\begin{iband}[signal]
Nothing in this module makes the model more capable. Everything in it is about \textbf{what the model is connected to}, and about \hilite{what happens when it is wrong}
\end{iband}
\end{frame}

\framekicker{Vocabulary}
\begin{frame}[shrink=25]{Five terms, in increasing order of autonomy}
\begin{center}
\includegraphics[width=0.92\linewidth,height=0.60\textheight,keepaspectratio]{figs/hendri-ladder.tex.pdf}
\end{center}
\vskip 3pt{\color{ink3}\fontsize{7.4}{9.4}\selectfont Each step gives the model more say over what happens next.\par}
{\fontsize{9.4}{12.6}\selectfont\color{fg2}These are used interchangeably in vendor material and they are not interchangeable. \textbf{The difference between them is a cost difference, a testability difference, and a risk difference} --- which makes it worth being precise.\par}\vskip 4pt
\end{frame}

\framekicker{The operational test}
\begin{frame}[shrink=25]{One question separates a workflow from an agent}
\begin{center}
\includegraphics[width=0.92\linewidth,height=0.60\textheight,keepaspectratio]{figs/hendri-who-decides.tex.pdf}
\end{center}
\vskip 3pt{\color{ink3}\fontsize{7.4}{9.4}\selectfont Everything else --- the model, the tools, the prompt --- can be identical.\par}
\begin{iband}[signal]
\textbf{If you can draw the flowchart in advance and it does not change per request, you have a workflow.} Build it as a workflow. It will be cheaper, faster, and you will be able to test it.
\end{iband}
\end{frame}
\note{Press on this. Most systems described as agents in industry talks are workflows with a model in one or two of the boxes, and that is a compliment, not a criticism.}

\framekicker{Vocabulary}
\begin{frame}[shrink=25]{The five, stated precisely}
\begin{itable}{>{\raggedright\arraybackslash}p{0.2068\textwidth}>{\raggedright\arraybackslash}p{0.3572\textwidth}>{\raggedright\arraybackslash}p{0.376\textwidth}}
\thead{Term} & \thead{Control flow} & \thead{What it is good for} \\ \midrule
\textbf{AI} & None --- one call in, one answer out & Classification, extraction, summarization, generation. \\
\textbf{AI workflow} & \textbf{Yours}, fixed at design time & Most business processes. Predictable cost and latency. \\
\textbf{AI agent} & \textbf{The model's}, chosen per step from a tool set & Tasks whose shape is not known until you see the request. \\
\textbf{Agentic system} & The model's, plus memory and planning across turns & Long-running work; tasks that resume. \\
\textbf{Multi-agent} & A supervisor routing between specialists & Genuinely separable subtasks with different tools or permissions. \\
\end{itable}
\end{frame}

\framekicker{A note on the hype}
\begin{frame}[shrink=25]{Deflating this on purpose}
\leadpar{\textbf{Most problems brought to "agents" are workflows.} Most multi-agent designs do not survive contact with a cost model.}
{\fontsize{9.4}{12.6}\selectfont\color{fg2}That is not a reason to avoid the topic. It is the reason to learn it properly --- because the value in this area is almost entirely in \textbf{tools, memory, and evaluation}, and almost none of it is in prompts.\par}\vskip 4pt
{\fontsize{9.4}{12.6}\selectfont\color{fg2}A useful habit for the rest of this module: every time you see a technique, ask \textbf{what it would cost to run 10,000 times a day, and how you would know it was working}.\par}\vskip 4pt
\end{frame}

\coursesection{02}{How an agent actually runs}{A loop, a tool schema, and something that makes it stop.}

\framekicker{The loop}
\begin{frame}[shrink=25]{Plan, act, observe, check}
\begin{center}
\includegraphics[width=0.92\linewidth,height=0.60\textheight,keepaspectratio]{figs/hendri-loop.tex.pdf}
\end{center}
\vskip 3pt{\color{ink3}\fontsize{7.4}{9.4}\selectfont The whole of agent engineering is in what fills each of these boxes.\par}
{\fontsize{9.4}{12.6}\selectfont\color{fg2}Note the shape. This is \textbf{the generation loop from chapter 16 with a tool call in the middle} --- the model's output at one step becomes part of its input at the next. Everything chapter 15 said about that pattern still applies.\par}\vskip 4pt
\end{frame}

\framekicker{The loop $\cdot$ the part people forget}
\begin{frame}[shrink=25]{What terminates it}
{\fontsize{9.4}{12.6}\selectfont\color{fg2}A loop with no exit condition is a bug that bills by the token. Every agent needs \textbf{all} of the following, not one of them.\par}\vskip 4pt
\begin{ibullets}
\item \textbf{Goal satisfied} --- the model says it is done, and something else checks that claim.
\item \textbf{Step budget} --- a hard maximum number of iterations.
\item \textbf{Token or currency budget} --- a hard maximum spend per request.
\item \textbf{Wall-clock timeout} --- because a hung tool is not the model's problem to solve.
\item \textbf{Repetition detector} --- the same tool with the same arguments twice in a row is almost always a loop, not progress.
\item \textbf{Escalation path} --- what happens when the budget is spent and the goal is not met. \textit{Silently returning a partial answer is the worst option.}
\end{ibullets}
\end{frame}

\framekicker{Tools}
\begin{frame}[shrink=25]{The model never touches anything}
\begin{center}
\includegraphics[width=0.92\linewidth,height=0.60\textheight,keepaspectratio]{figs/hendri-tool.tex.pdf}
\end{center}
\vskip 3pt{\color{ink3}\fontsize{7.4}{9.4}\selectfont The model emits a structured request. Your code decides whether to honour it.\par}
\begin{iband}[signal]
This is the single most important architectural fact in the module. \textbf{The model produces text describing an intention. Your code is the only thing that executes.} Every permission boundary you have lives in that gap.
\end{iband}
\end{frame}

\framekicker{Tools $\cdot$ design}
\begin{frame}[shrink=25]{A tool is an API designed for a reader who forgets}
\begin{isteps}
\item \textbf{Name it for the task, not the system.} \inlinecode{find\_customer\_by\_email} beats \inlinecode{crm\_query\_v2}.
\item \textbf{Write the description for the model, not for your team.} It is the only documentation the model gets, and it is read fresh every single call.
\item \textbf{Type every parameter, and validate on arrival.} Assume the arguments are adversarial, because sometimes they will be.
\item \textbf{Return errors as data, not exceptions.} A tool that returns \inlinecode{\{"error": "no customer with that email"\}} lets the agent recover; one that throws ends the run.
\item \textbf{Keep results small.} Every byte a tool returns occupies context that the model then pays to read on every subsequent step.
\end{isteps}
{\fontsize{9.4}{12.6}\selectfont\color{fg2}Most \textit{prompt engineering} problems in agent systems are actually \textbf{tool description problems} wearing a disguise.\par}\vskip 4pt
\end{frame}

\framekicker{Tools $\cdot$ the dangerous ones}
\begin{frame}[shrink=25]{Read tools and write tools are not the same risk}
\begin{minipage}[t]{0.494\textwidth}
\begin{icardgood}
\cardh{Read tools}
Search, retrieve, look up, calculate. A wrong call wastes tokens and returns something unhelpful. \textbf{Recoverable within the loop.}
\end{icardgood}
\end{minipage}%
\hfill
\begin{minipage}[t]{0.494\textwidth}
\begin{icardbad}
\cardh{Write tools}
Send, delete, transfer, approve, publish. A wrong call \textbf{changes the world}, and no amount of subsequent reasoning undoes it.
\end{icardbad}
\end{minipage}%
\begin{iband}[amber]
Treat these differently from day one. \textbf{Read tools can be autonomous; write tools should start behind an approval gate} and only move out of it with evidence from your evaluation harness.
\end{iband}
\end{frame}

\framekicker{Memory}
\begin{frame}[shrink=25]{Four kinds, solving four different problems}
\begin{center}
\includegraphics[width=0.92\linewidth,height=0.60\textheight,keepaspectratio]{figs/hendri-memory.tex.pdf}
\end{center}
\vskip 3pt{\color{ink3}\fontsize{7.4}{9.4}\selectfont Borrowed vocabulary from cognitive psychology, used loosely --- but the four categories map onto four different pieces of infrastructure.\par}
{\fontsize{9.4}{12.6}\selectfont\color{fg2}The mistake to avoid is treating \textit{memory} as one thing and reaching for a vector database by reflex. \textbf{Most agents need working memory managed well and nothing else.}\par}\vskip 4pt
\end{frame}

\framekicker{Memory $\cdot$ the practical version}
\begin{frame}[shrink=25]{What each one costs you}
\begin{itable}{>{\raggedright\arraybackslash}p{0.188\textwidth}>{\raggedright\arraybackslash}p{0.3572\textwidth}>{\raggedright\arraybackslash}p{0.3948\textwidth}}
\thead{Kind} & \thead{Implementation} & \thead{The failure mode it introduces} \\ \midrule
\textbf{Working} & The context window itself & Overflow. Old steps silently fall out and the agent forgets its goal. \\
\textbf{Episodic} & A transcript store, keyed by session & Replaying stale state as if it were current. \\
\textbf{Semantic} & Vector or keyword index over documents --- chapter 16's RAG & Retrieving confidently irrelevant context, which is worse than none. \\
\textbf{Procedural} & Versioned prompts, tool sets, playbooks & Drift between what is deployed and what was evaluated. \\
\end{itable}
{\fontsize{9.4}{12.6}\selectfont\color{fg2}Each row is a system to operate. \textbf{Add them one at a time, and only when a specific observed failure demands it.}\par}\vskip 4pt
\end{frame}

\coursesection{03}{The tech stack}{Six layers, and the two that teams skip.}

\framekicker{Architecture}
\begin{frame}[shrink=25]{Six layers}
\begin{center}
\includegraphics[width=0.92\linewidth,height=0.60\textheight,keepaspectratio]{figs/hendri-stack.tex.pdf}
\end{center}
\vskip 3pt{\color{ink3}\fontsize{7.4}{9.4}\selectfont Guardrails and observability are the two layers most first attempts leave out entirely.\par}
\end{frame}

\framekicker{Architecture $\cdot$ model layer}
\begin{frame}[shrink=25]{Hosted or self-hosted, and what actually decides it}
\begin{minipage}[t]{0.494\textwidth}
\begin{icard}
\cardh{Hosted API}
Best capability per unit of effort, no infrastructure, immediate access to new models. \textbf{Cost scales linearly with usage}, and your data leaves your perimeter.
\end{icard}
\end{minipage}%
\hfill
\begin{minipage}[t]{0.494\textwidth}
\begin{icard}
\cardh{Self-hosted open weights}
Data stays inside. Fixed infrastructure cost rather than per-token. \textbf{Lower capability at the same price point}, and you now operate a GPU fleet.
\end{icard}
\end{minipage}%
\begin{iband}[amber]
In practice the decision is rarely about capability or cost. It is about \textbf{where the data is allowed to be} --- which is a regulatory question, not an engineering one, and it should be answered before any code is written.
\end{iband}
\end{frame}
\note{For participants from regulated industries this is the slide that matters most. The constraint is usually a specific clause in a specific rule, and the right move is to go and read it rather than assume.}

\framekicker{Architecture $\cdot$ orchestration layer}
\begin{frame}[shrink=25]{What the orchestrator is responsible for}
\begin{ibullets}
\item \textbf{Running the loop} --- and enforcing every one of the termination conditions from earlier.
\item \textbf{Routing tool calls} to implementations, and validating arguments before they arrive.
\item \textbf{Retrying} transient failures, with backoff, without re-charging the model for the same reasoning.
\item \textbf{Accounting} --- tokens, currency, and latency, per request and per tool.
\item \textbf{Emitting traces} for everything above.
\end{ibullets}
{\fontsize{9.4}{12.6}\selectfont\color{fg2}Frameworks exist for all of this and are worth using. But note what the list contains: \textbf{queueing, retries, budgets, and telemetry.} This is ordinary distributed-systems work, and your existing platform team already knows how to do it.\par}\vskip 4pt
\end{frame}

\framekicker{Architecture $\cdot$ guardrails}
\begin{frame}[shrink=25]{Three places to put a check}
\begin{isteps}
\item \textbf{On the input.} Reject or sanitise before the model sees it --- prompt injection carried in a retrieved document is the case people miss.
\item \textbf{On the tool call.} Validate arguments, check permissions against the \textit{end user's} identity rather than the service's, and gate write tools behind approval.
\item \textbf{On the output.} Schema-check structured output; scan free text for the categories your policy prohibits; verify claims against retrieved sources where you can.
\end{isteps}
\begin{iband}[signal]
A guardrail that is itself a model call is a guardrail that can be wrong. \textbf{Prefer deterministic checks wherever the check can be expressed deterministically} --- schema validation, allow-lists, permission lookups.
\end{iband}
\end{frame}

\framekicker{Architecture $\cdot$ the one people skip}
\begin{frame}[shrink=25]{Observability is not optional here}
{\fontsize{9.4}{12.6}\selectfont\color{fg2}In a deterministic system, a bug reproduces. In this one it does not --- the same request can take a different path on Tuesday than it did on Monday, with the same code and the same weights.\par}\vskip 4pt
\leadpar{\textbf{A trace of every step, every tool call, and every token is the only debugging tool you have.} Build it first, not after the first incident.}
{\fontsize{9.4}{12.6}\selectfont\color{fg2}The minimum useful trace records: the full prompt sent, the model and version, every tool call with arguments and results, token counts and cost per step, latency per step, and the termination reason. \hilite{If you cannot answer \textit{why did it do that} from your logs, you do not have logs.}\par}\vskip 4pt
\end{frame}

\coursesection{04}{Multi-agent, and when not to}{A supervisor, some specialists, and a three-question test.}

\framekicker{Multi-agent}
\begin{frame}[shrink=25]{The supervisor pattern}
\begin{center}
\includegraphics[width=0.92\linewidth,height=0.60\textheight,keepaspectratio]{figs/hendri-multiagent.tex.pdf}
\end{center}
\vskip 3pt{\color{ink3}\fontsize{7.4}{9.4}\selectfont Each specialist has its own tools, its own context, and its own instructions.\par}
{\fontsize{9.4}{12.6}\selectfont\color{fg2}This is the pattern that appears in almost every multi-agent framework, under various names. The variations --- peer-to-peer messaging, hierarchies, debate --- are \textbf{elaborations of it, and each elaboration multiplies the cost.}\par}\vskip 4pt
\end{frame}

\framekicker{Multi-agent $\cdot$ the test}
\begin{frame}[shrink=25]{Three questions before you split}
\begin{center}
\includegraphics[width=0.92\linewidth,height=0.60\textheight,keepaspectratio]{figs/hendri-when-multi.tex.pdf}
\end{center}
\vskip 3pt{\color{ink3}\fontsize{7.4}{9.4}\selectfont If none of the three is a yes, you want one agent with more tools.\par}
\begin{iband}[signal]
\textbf{Splitting for conceptual tidiness is the common mistake.} Roles named \textit{Researcher}, \textit{Writer}, and \textit{Critic} feel organised and typically cost three times as much as one agent doing the same work with the same tools.
\end{iband}
\end{frame}

\framekicker{Multi-agent $\cdot$ the honest arithmetic}
\begin{frame}[shrink=25]{What each additional agent costs}
\begin{itable}{>{\raggedright\arraybackslash}p{0.2444\textwidth}>{\raggedright\arraybackslash}p{0.6956\textwidth}}
\thead{Cost} & \thead{Why it grows} \\ \midrule
\textbf{Tokens} & Every hand-off re-states context. A supervisor that summarises for three specialists pays for that summary three times, plus its own reasoning. \\
\textbf{Latency} & Sequential hand-offs add up. Parallel ones only help if the subtasks are genuinely independent. \\
\textbf{Failure surface} & Each agent can loop, drift, or hallucinate independently --- \textbf{and one specialist's wrong answer becomes another's trusted input.} \\
\textbf{Evaluation} & You now need per-agent evals \textit{and} end-to-end evals, because a system that is right overall can be right for the wrong reasons. \\
\end{itable}
{\fontsize{9.4}{12.6}\selectfont\color{fg2}A finding worth carrying: decompose a task into five or six specialists and it is common for \textbf{only one or two to justify their existence} once the traces are read.\par}\vskip 4pt
\end{frame}
\note{This is the slide to be blunt on. The room will have seen impressive multi-agent diagrams; almost none of them come with a cost table.}

\framekicker{Multi-agent $\cdot$ when it does pay}
\begin{frame}[shrink=25]{Three shapes where splitting is right}
\begin{minipage}[t]{0.3253\textwidth}
\begin{icardgood}
\cardh{Different permissions}
One agent may read customer records; another may write to the ledger. \textbf{Splitting is how you keep those capabilities apart} --- a security boundary, not a stylistic one.
\end{icardgood}
\end{minipage}%
\hfill
\begin{minipage}[t]{0.3253\textwidth}
\begin{icardgood}
\cardh{Genuine parallelism}
Twenty documents to analyse independently. Fan out, collect, assemble. The latency win is real and the contexts do not interact.
\end{icardgood}
\end{minipage}%
\hfill
\begin{minipage}[t]{0.3253\textwidth}
\begin{icardgood}
\cardh{Context that will not fit}
A task whose inputs genuinely exceed the window. Splitting is a \textbf{memory strategy} here, and the supervisor's summaries are the compression.
\end{icardgood}
\end{minipage}%
\end{frame}

\coursesection{05}{Evaluation and operation}{How you know it works, and what you do when it stops.}

\framekicker{Evaluation}
\begin{frame}[shrink=25]{You cannot test this the way you test code}
{\fontsize{9.4}{12.6}\selectfont\color{fg2}The system is non-deterministic, the output space is open, and the correct answer is often a judgement call. Standard unit testing does not reach any of that.\par}\vskip 4pt
\begin{center}
\includegraphics[width=0.92\linewidth,height=0.60\textheight,keepaspectratio]{figs/hendri-eval.tex.pdf}
\end{center}
\vskip 3pt{\color{ink3}\fontsize{7.4}{9.4}\selectfont Traces first. Everything downstream is built on them.\par}
\end{frame}

\framekicker{Evaluation $\cdot$ the three levels}
\begin{frame}[shrink=25]{What each level catches}
\begin{itable}{>{\raggedright\arraybackslash}p{0.2256\textwidth}>{\raggedright\arraybackslash}p{0.376\textwidth}>{\raggedright\arraybackslash}p{0.3384\textwidth}}
\thead{Level} & \thead{What it asserts} & \thead{What it misses} \\ \midrule
\textbf{Unit evals} & Given this input, the right tool is called with the right arguments. & Whether the sequence of correct calls produces a useful outcome. \\
\textbf{End-to-end evals} & Given this task, the final outcome matches a known good one. & \textit{Why} it succeeded --- a right answer for the wrong reason still passes. \\
\textbf{Human review} & Whether the output is actually good, by a standard you cannot encode. & Coverage. It is sampled, so it finds classes of failure, not instances. \\
\end{itable}
\begin{iband}[signal]
Run all three \textbf{before every prompt change, every model upgrade, and every tool change}. In this domain a prompt edit is a code change with no type checker behind it.
\end{iband}
\end{frame}

\framekicker{Evaluation $\cdot$ building the set}
\begin{frame}[shrink=25]{Where the evaluation tasks come from}
\begin{isteps}
\item \textbf{Start with twenty real requests}, taken from whatever the process looks like today. Not twenty imagined ones.
\item \textbf{Write down the good outcome for each} --- before you build anything, while you are still honest about what good means.
\item \textbf{Add every failure you observe} in development, permanently. This is how the set grows, and it is why it is worth versioning.
\item \textbf{Include the boring cases.} A set made only of hard cases will let you ship something that fails the easy ones.
\item \textbf{Keep a held-out slice} you do not look at while iterating. Chapter 18's warning about validation-set overfitting applies here exactly.
\end{isteps}
\end{frame}

\framekicker{Operation}
\begin{frame}[shrink=25]{Four failure modes, and how each announces itself}
\begin{center}
\includegraphics[width=0.92\linewidth,height=0.60\textheight,keepaspectratio]{figs/hendri-failure.tex.pdf}
\end{center}
\vskip 3pt{\color{ink3}\fontsize{7.4}{9.4}\selectfont None of these throw an exception. All of them show in a trace.\par}
{\fontsize{9.4}{12.6}\selectfont\color{fg2}Set alerts on the \textbf{shape} of runs, not only on errors: mean steps per request, cost per request, tool-call repetition rate, and escalation rate. \textbf{A rising step count is usually the first visible sign that something upstream changed.}\par}\vskip 4pt
\end{frame}

\framekicker{Operation $\cdot$ deployment}
\begin{frame}[shrink=25]{The autonomy ladder}
\begin{center}
\includegraphics[width=0.92\linewidth,height=0.60\textheight,keepaspectratio]{figs/hendri-maturity.tex.pdf}
\end{center}
\vskip 3pt{\color{ink3}\fontsize{7.4}{9.4}\selectfont Move up a rung only with evidence from the evaluation harness.\par}
\begin{iband}[signal]
At \textbf{L3}, the question that matters is not \textit{how good is it?} but \hilite{what is the blast radius when it is wrong?} Define that boundary in code --- spend limits, record scopes, reversibility --- not in the prompt.
\end{iband}
\end{frame}

\coursesection{06}{Cases}{What we will build, and what each one is chosen to teach.}

\framekicker{Cases}
\begin{frame}[shrink=25]{Single-agent cases}
\begin{itable}{>{\raggedright\arraybackslash}p{0.3008\textwidth}>{\raggedright\arraybackslash}p{0.6392\textwidth}}
\thead{Case} & \thead{What it teaches} \\ \midrule
\textbf{Document Q\&A with citations} & Retrieval as a tool, not as a preprocessing step. Grounding, and how to make a citation verifiable rather than decorative. \\
\textbf{Structured extraction from messy input} & Schema-constrained output, validation on arrival, and what to do with the 20\% that fails validation. \\
\textbf{Internal API assistant} & Tool design under real permissions. Read tools autonomous, write tools gated. Identity propagation. \\
\textbf{Data analysis with code execution} & The most powerful and most dangerous tool there is. Sandboxing, resource limits, and why the output must still be checked. \\
\end{itable}
\end{frame}

\framekicker{Cases}
\begin{frame}[shrink=25]{Multi-agent cases}
\begin{itable}{>{\raggedright\arraybackslash}p{0.3008\textwidth}>{\raggedright\arraybackslash}p{0.6392\textwidth}}
\thead{Case} & \thead{What it teaches} \\ \midrule
\textbf{Parallel document triage} & The clean case for fan-out: independent subtasks, real latency win, no shared context. The one that unambiguously justifies splitting. \\
\textbf{Research and draft with a separate reviewer} & The \textit{tidy} decomposition --- and a cost comparison against one agent doing both, run honestly. \\
\textbf{Cross-permission workflow} & Two agents that must not share capabilities. Splitting as a security boundary. \\
\textbf{Escalation to a human} & The hand-off nobody designs and everybody needs. What state goes with it, and how the human hands it back. \\
\end{itable}
{\fontsize{9.4}{12.6}\selectfont\color{fg2}All four ship in the \textbf{ai-agentic-demo} monorepo, each with a trace viewer and an evaluation set, so the cost and quality claims on these slides can be checked rather than believed.\par}\vskip 4pt
\end{frame}

\framekicker{Cases $\cdot$ the exercise}
\begin{frame}[shrink=25]{What you will do with them}
\begin{isteps}
\item \textbf{Bring a process from your own work} --- one that exists, with real inputs you can describe.
\item \textbf{Classify it} using the operational test from section 01. Most will be workflows, and that is the correct answer.
\item \textbf{Write the twenty evaluation cases} before writing any prompt.
\item \textbf{Build the smallest version} that could work: one agent, the fewest tools that cover the task.
\item \textbf{Run it, read the traces, and report what surprised you.} That last step is the assessment.
\end{isteps}
\begin{iband}[signal]
The deliverable is not a working agent. It is \textbf{a defensible answer to whether you should have built one} --- supported by traces, costs, and an evaluation set.
\end{iband}
\end{frame}

\framekicker{Summary}
\begin{frame}[shrink=25]{What has to survive this module}
\begin{isteps}
\item \textbf{Who decides the next step} is the question that separates a workflow from an agent. Most things are workflows; build them as workflows.
\item \textbf{The loop needs more than one exit.} Goal, steps, tokens, time, repetition, and an escalation path.
\item \textbf{The model never touches anything.} It emits a structured intention; your code decides whether to honour it. Every permission boundary lives there.
\item \textbf{Read tools and write tools are different risks}, and should have different defaults from day one.
\item \textbf{Memory is four things, not one.} Add each only when an observed failure demands it.
\end{isteps}
\end{frame}

\framekicker{Summary}
\begin{frame}[shrink=25]{What has to survive this module (2 of 2)}
\begin{isteps}
\item \textbf{Guardrails and observability are layers}, not features. Traces come before everything else, because nothing here reproduces.
\item \textbf{Three questions before splitting into multiple agents} --- different permissions, real parallelism, or context that will not fit. Tidiness is not one of them.
\item \textbf{Evaluate at three levels} --- unit, end-to-end, and sampled human review --- and run all three before any prompt or model change.
\item \textbf{Move up the autonomy ladder on evidence}, and define the blast radius in code rather than in the prompt.
\end{isteps}
\vskip 4pt
\reslink{REPO}{ai-agentic-demo}{https://github.com/situkangsayur/ai-agentic-demo}
\reslink{BACK}{Chapter 16 --- Text generation}{../ch16/index.html}
\vskip 2pt
\end{frame}

\end{document}
